\documentclass[12pt]{article}
\usepackage[utf8]{inputenc}
\usepackage{amsmath, amssymb, amsfonts, mathtools, array}
\usepackage{geometry}
\geometry{a4paper, margin=1in}
\title{Student Answer Script Transcription (High-Resolution Stitched Document)}
\author{Roll Number: 210103096}
\date{\today}
\begin{document}
\maketitle

% ==================================================
% Content from Slice 1: output\_qwen\_vl\_page-210103096-01.tex
% ==================================================
\section*{Solution}

\subsection*{Question 1}

Given,
\_\_MATH\_BLOCK\_0\_\_
   \_\_MATH\_BLOCK\_1\_\_
   Since \(\_\_MATH\_BLOCK\_2\_\_\) is not in S, the zero matrix is not in S.

2. \textbf{If \( x \in S \) and \( y \in S \), then \( x + y \in S \):}
   Let \( x = \_\_MATH\_BLOCK\_3\_\_ \) and \( y = \_\_MATH\_BLOCK\_4\_\_ \)
   \_\_MATH\_BLOCK\_5\_\_
   Since \(-0.5\) is not in S, \(x + y \neq S\).

3. \textbf{If \( x \in S \) and \( y \in S \), then \( xy \in S \):}
   Let \( x = \_\_MATH\_BLOCK\_6\_\_ \) and \( y = \_\_MATH\_BLOCK\_7\_\_ \)
   \_\_MATH\_BLOCK\_8\_\_
   Since \( xy \notin S \), condition 3 is not satisfied.

Therefore, \( S \) is not a subspace of \( M\_n(\mathbb{R}) \).

% ==================================================
% Content from Slice 2: output\_qwen\_vl\_page-210103096-02.tex
% ==================================================
\section*{Solution}

\subsection*{Question 1}

The given matrix is:

\_\_MATH\_BLOCK\_9\_\_

To find the basis for the null space of A, we need to perform row reduction on the matrix.

Step-by-step explanation:

1. \textbf{Row Reduction:}
   - Subtract \( \frac{3}{7} \) times row 1 from row 2.
   - Subtract \( \frac{9}{7} \) times row 1 from row 3.
   - Add \( \frac{6}{7} \) times row 1 to row 4.

This results in:

\_\_MATH\_BLOCK\_10\_\_

2. \textbf{Identify the Basis:}
   - The resulting matrix has only one non-zero row, which means that the null space of A is spanned by the first column vector.

Therefore, the basis for the null space of A is:

\_\_MATH\_BLOCK\_11\_\_

So, the answer is:

\_\_MATH\_BLOCK\_12\_\_

\end{document}
